% Part: second-order-logic
% Chapter: sol-and-sets
% Section: cardinalities

\documentclass[../../../include/open-logic-section]{subfiles}

\begin{document}

\olfileid{sol}{set}{crd}

\olsection{संचांची आकारमाने}

\begin{explain}
प्रांत सांत किंवा अनंत, !!{enumerable} किंवा
!!{nonenumerable} आहे हे जसे व्यक्त करता येते,
तसेच~$\Domain{M}$ चा उपसंच सांत किंवा अनंत,
!!{enumerable} किंवा !!{nonenumerable} असण्याचा
गुणधर्मही परिभाषित करता येतो.
\end{explain}

\begin{prop}
%READERNOTE{SOL6-A097: मूळ Inf, Count आणि Aleph-one सूत्रांचे आधी निदर्शित दोष येथे प्रत्यक्ष दुरुस्त केले आहेत. X-to-X closure, restricted injectivity, empty-set countability आणि strictly smaller cardinality च्या उपसंचांची अट आवश्यक आहेत.}
सूत्र $\fn{Inf}(X) \ident$
\begin{multline*}
\lexists[u][(\lforall[x][(X(x) \lif X(u(x)))] \land {} \\
\lforall[x][\lforall[y][((X(x) \land X(y) \land \eq[u(x)][u(y)]) \lif
      \eq[x][y])]] \land {} \\
\lexists[y][(X(y) \land \lforall[x][(X(x) \lif \eq/[y][u(x)])])])]
\end{multline*}
याची चल-मूल्यांकन~$s$ नुसार पूर्ति होईल, जर आणि
फक्त जर $s(X)$ अनंत असेल. स्रोतदुरुस्ती: फलन
X ला स्वतःमध्येच पाठवते, X वर एकास एक असते आणि X मधील
किमान एक घटक त्याच्या X वरील व्याप्तीत नसतो, अशा सर्व अटी
घेतल्या आहेत. त्यामुळे मूळ सूत्राचा सांत उपसंचावरील counterexample
दूर होतो. अनंततेचा नेहमीचा निवडीसह संदर्भ येथेही लागू आहे.
\end{prop}

\begin{prop}
सूत्र $\fn{Count}(X) \ident $
\begin{multline*}
\lnot\lexists[z][X(z)] \lor {} \\
\lexists[z][\lexists[u][(X(z) \land
    \lforall[x][(X(x) \lif X(u(x)))] \land {} \\
\lforall[Y][((Y \subseteq X \land Y(z) \land
      \lforall[x][(Y(x) \lif Y(u(x)))]) \lif X = Y)])]]
\end{multline*}
याची चल-मूल्यांकन~$s$ नुसार पूर्ति होईल, जर आणि
फक्त जर $s(X)$ !!{enumerable} असेल. स्रोतदुरुस्ती:
पहिला खंड रिकामा संच घेतो. अरिक्त X साठी आरंभीचा घटक
आणि X ला स्वतःमध्ये पाठवणारे फलन घेतले आहे. शेवटची अट
फक्त X मधील बंद उपसंचांबद्दल आहे; त्यामुळे X हा आरंभीच्या
घटकाचा पूर्ण orbit असतो. हे संपूर्ण प्रांतच X असण्याची
चुकीची सक्ती करत नाही.
\end{prop}

कँटरच्या प्रमेयावरून !!{nonenumerable} संच आहेत,
आणि अनंत आकारमानांचे निरनिराळे अमर्याद स्तरही
आहेत, हे आपल्याला माहीत आहे. संचसिद्धान्तात
संचांच्या आकारमानांचे संपूर्ण अंकगणित विकसित केले
जाते आणि संचांना अनंत संचांक दिले जातात. सांत
संचांची आकारमाने मोजणारे संचांक म्हणजे स्वाभाविक
संख्या. !!{denumerable} संचांचा संचांक हा पहिला
अनंत संचांक असून त्याला~$\aleph_0$ (``अलेफ-शून्य'')
म्हणतात. पुढील अनंत आकारमान~$\aleph_1$ आहे.
गणनीय नसलेल्या संचाला असू शकणारे हे
सर्वांत लहान आकारमान आहे (गणनीय अनंत संचाचे
आकारमान~$\aleph_0$ असते). ``$X$ चे आकारमान
$\aleph_0$ आहे'' याची व्याख्या
$\fn{Aleph}_0(X) \liff \fn{Inf}(X) \land \fn{Count}(X)$
अशी करता येते. $X$ चे आकारमान $\aleph_1$
असेल, जर आणि फक्त जर त्याच्यापेक्षा काटेकोरपणे लहान
आकारमानाचे सर्व उपसंच सांत किंवा~$\aleph_0$ आकारमानाचे
असतील आणि तो स्वतः सांत किंवा~$\aleph_0$ आकारमानाचा
नसेल. म्हणून पुढील दुरुस्त सूत्र घेतो:
$\fn{Aleph_1}(X) \ident \lnot\fn{Count}(X) \land
\lforall[Y][((Y \subseteq X \land \lnot\cardeq{Y}{X}) \lif
  (\lnot\fn{Inf}(Y) \lor \fn{Aleph}_0(Y)))]$. $\aleph_2$ आकारमानाची
व्याख्याही अशाच प्रकारे करायची, असे त्यात म्हटले आहे.
स्रोतदुरुस्ती: X अगणनीय असावा लागतो. स्वतः X किंवा
त्याच्याशी तुल्यबल असलेला उपसंच हा लहान आकारमानाचा नाही;
त्याला finite किंवा aleph-zero असण्याची चुकीची अट लावलेली नाही.

एक विशेष महत्त्वाचे आकारमान म्हणजे तथाकथित
सांतत्याचा संचांक. ते $\Pow{\Nat}$ चे आकारमान,
किंवा समतुल्यपणे~$\Real$ चे आकारमान आहे.
एखादा संच सांतत्याच्या आकारमानाचा आहे, हेही
द्वितीय-क्रम तर्कशास्त्रात व्यक्त करता येते; मात्र
त्यासाठी आणखी थोडे काम करावे लागते.

\end{document}
